\subsection*{第3章の索引用語}

\begin{description}
  \item[\term{OS}] Operating Systemの略で、計算資源を分配し共通の実行環境を提供します。
  \item[\term{ユーザー空間}] 通常のアプリケーションが制限された権限で動く実行領域です。
  \item[\term{コンテキストスイッチ}] CPUが実行対象を切り替えるために状態を保存し復元する処理です。
  \item[\term{スケジューラ}] 実行可能な処理へCPU時間を割り当てるOSの機能です。
  \item[\term{マルチプロセス}] 複数のプロセスへ仕事を分け、原則としてアドレス空間を分離する構成です。
  \item[\term{マルチスレッド}] 一つのプロセスに複数の実行の流れを置き、ヒープなどを共有する構成です。
  \item[\term{並行}] 複数の仕事が開始から終了までの時間を重ねて進む性質です。
  \item[\term{並列}] 複数の仕事が同じ時点に物理的に実行される性質です。
  \item[\term{ガベージコレクション}] 到達不能になったヒープ領域などを処理系が自動で回収する仕組みです。
  \item[\term{ファイルディスクリプタ}] プロセスが開いたファイルやソケットを参照する小さな整数です。
  \item[\term{標準入力}] プロセスが既定で入力を受け取るストリームです。
  \item[\term{標準出力}] プロセスが既定で結果を書き出すストリームです。
  \item[\term{ページテーブル}] 仮想ページと物理メモリの対応を保持する表です。
  \item[\term{TLB}] Translation Lookaside Bufferの略で、アドレス変換結果を保持するキャッシュです。
  \item[\term{int}] Goで符号付き整数を表す基本型です。
  \item[\term{string}] Goで読み取り専用のバイト列として文字列を表す型です。
  \item[\term{bool}] Goで真偽値を表す型です。
  \item[\term{slice}] Goで配列の区間を参照する記述子を持つ型です。
  \item[\term{map}] Goでキーと値の対応を保持する型です。
  \item[\term{struct}] Goで複数のフィールドを一つの値へまとめる型です。
  \item[\term{サイズ}] 値やデータ構造が占めるバイト数です。
  \item[\term{参照の共有}] 複数の値が同じ可変データを指す状態です。
  \item[\term{連続メモリ}] 要素がアドレス上で隣接して置かれるメモリ領域です。
  \item[\term{インデックス}] 列の中の要素位置を指定する整数です。
  \item[\term{スライスヘッダ}] 元配列への参照、長さ、容量を持つスライスの記述子です。
  \item[\term{len}] Goで値の長さを返す組み込み関数です。
  \item[\term{cap}] Goで値の容量を返す組み込み関数です。
  \item[\term{append}] Goでスライスへ要素を加えた新しいスライス値を返す組み込み関数です。
  \item[\term{元配列}] スライスが要素を参照する背後の配列です。
  \item[\term{ハッシュテーブル}] ハッシュ値を用いてキーから格納位置を探すデータ構造です。
  \item[\term{バケット}] ハッシュテーブルで複数の要素候補をまとめて置く領域です。
  \item[\term{Pythonオブジェクト}] Pythonでidentity、type、valueを持つ値の実体です。
  \item[\term{identity}] Pythonオブジェクトを生存期間中に一意に識別する性質です。
  \item[\term{type}] Pythonオブジェクトが可能な操作を定める種類です。
  \item[\term{value}] Pythonオブジェクトが表す内容です。
  \item[\term{ミュータブル}] 作成後に内容を変更できる性質です。
  \item[\term{イミュータブル}] 作成後に同じオブジェクトの内容を変更しない性質です。
\end{description}
